\chapter{\IfLanguageName{dutch}{Stand van zaken}{State of the art}}%
\label{ch:stand-van-zaken}

% Tip: Begin elk hoofdstuk met een paragraaf inleiding die beschrijft hoe
% dit hoofdstuk past binnen het geheel van de bachelorproef. Geef in het
% bijzonder aan wat de link is met het vorige en volgende hoofdstuk.

% Pas na deze inleidende paragraaf komt de eerste sectiehoofding.


In dit hoofdstuk wordt de theoretische en contextuele achtergrond gepresenteerd die nodig is voor het bestuderen van transparantie en herleidbaarheid van impactanalyses bij VLAIO. Eerst worden de begrippen data lineage en provenance toegelicht, gevolgd door de rol van lineage binnen data governance en impactanalyse. Daarna wordt besproken hoe de meerwaarde van lineage in een organisatiecontext kan worden geëvalueerd. Daarna worden metadata management, data catalogs en de positionering van OpenMetadata besproken. Tot slot wordt de onderzoeksleemte afgebakend. Dit hoofdstuk vormt het referentiekader voor de interpretatie van de case study en de evaluaties die in de volgende hoofdstukken aan bod komen.\\

\section{Data lineage en provenance}

Bij VLAIO verwijst data lineage naar de volledige weg die data binnen een organisatie aflegt van de initiële bron tot het eindrapport of een andere analytische toepassing. Het beschrijft zowel de herkomst van gegevens, maar kan ook alle transformaties, uitbreidingen en afhankelijkheden van de data, doorheen verschillende systemen zichtbaar maken. Zoals \textcite{verma2024tracing} beschrijven, bewegen datasets met verschillende stappen door een data-pipeline, waarbij deze continu van vorm veranderen. Data lineage documenteert deze opeenvolgende stappen en geeft inzicht in hoe de data wordt verzameld, verwerkt en uiteindelijk gebruikt binnen de organisatie.\\

In de literatuur wordt naast de term data lineage ook de term data provenance gebruikt. Hoewel ze vaak door elkaar gebruikt worden \autocite{mayernik2013data}, hebben ze elk een specifiek doel \autocite{infobelpro2025}. Lineage richt zich op de geschiedenis van de data en dependency graphs tussen data-resources \autocite{mayernik2013data}, oftewel de documentatie hoe data verwerkt wordt en beweegt door systemen \autocite{patel2020data}. Provenance focust daarentegen op de oorsprong of bron van verwerkingsresultaten, de historische en contextuele achtergrond met informatie over systemen en processen \autocite{mayernik2013data,patel2020data}. In de praktijk worden beide idealiter gecombineerd. Op die manier ontstaat er een breder beeld van betrouwbaarheid, waarbij data lineage inzicht geeft in de dataflow en data provenance de oorsprong en authenticiteit van de data bewijst \autocite{infobelpro2025}.\\

Binnen data-intensieve omgevingen worden lineage en provenance gezien als cruciale componenten voor het garanderen van transparantie, herleidbaarheid en herbruikbaarheid van data \autocite{mayernik2013data}. Ze maken de verwerkingsstappen van de data inzichtelijk, wat zorgt voor meer vertrouwen en minder misinterpretaties bij gebruikers. Bovendien helpt dit om te beoordelen of de data van voldoende kwaliteit is en dus geschikt is voor specifieke toepassingen \autocite{mayernik2013data}. Ook het traceren van fouten wordt efficiënter. Dit is zeer relevant in complexe en gedistribueerde data-architecturen met meerdere databronnen, transformatielagen en systemen zoals bijvoorbeeld een cloudomgeving \autocite{patel2020data}.\\

Een gebrek aan een goed gestructureerde lineage of provenance kan leiden tot een verminderde transparantie bij de totstandkoming van bepaalde resultaten. Wanneer afhankelijkheden onvoldoende gedocumenteerd zijn, wordt het moeilijker om foutdetecties of impactanalyses uit te voeren, wanneer datasets of processen gewijzigd worden \autocite{patel2020data}. Hierdoor neemt dus de controleerbaarheid van rapportages en impactanalyses af, wat potentieel negatieve gevolgen heeft voor de besluitvorming binnen een organisatie \autocite{verma2024tracing}. \textcite{patel2020data} benadrukken bovendien dat onvolledige herleidbaarheid in gedistribueerde of cloud-omgevingen het identificeren van fouten aanzienlijk kan vertragen.\\

Een goede lineage-implementatie is volgens \textcite{tang2022modeling} niet zonder uitdagingen. Bestaande studies tonen aan dat implementaties vaak geen volledige SQL of procedurele talen ondersteunen, met als gevolg dat complexe SQL-queries en procedurele structuren dus moeilijk te traceren zijn. Daarnaast kan een gebrek aan multi-granulariteit op tuple/value-niveau en een beperkte ondersteuning voor complexe transformatieprocessen leiden tot een onvolledige afhankelijkheidsdetectie. Een effectieve toepassing vereist daarom robuuste modellen voor SQL-syntax en datatransformaties, naast governance voor consistente metadata.\\

\section{Data lineage binnen data governance}

Data governance omvat de principes en processen waarmee een organisatie haar data beheert, met aandacht voor beveiliging, regelgeving, rollen en verantwoordelijkheden, zodat data veilig en strategisch kan worden gebruikt \autocite{verma2024tracing}. \textcite{verma2024tracing} stellen dat een organisatie goede governance nodig heeft voor vlotte en duurzame operaties, met de focus op data-kwaliteit, veiligheid en compliance. Door het voorzien van gestandaardiseerde data, definities en validatie regels, zorgt governance voor een groter vertrouwen in de data binnen een organisatie.\\

Echter is het implementeren van governance geen eenvoudige taak en vereist meerdere stappen en pre-condities, zoals onder andere managementondersteuning executive support), duidelijke bedrijfsdoelstellingen en een volledige inventarisatie van de data \autocite{verma2024tracing}. Hoewel governance het vertrouwen in data via standaarden, definities en validatieregels wil verhogen, blijft dit moeilijk als er geen transparantie is in de oorsprong en de bewerkingen van de data. Dit is waar data lineage te pas komt. Het maakt zichtbaar waar de data vandaan komt, hoe ze wordt getransformeerd en waar ze naartoe gaat. \textcite{verma2024tracing} beschrijven dat data lineage een component van governance is die zorgt voor transparantie, de sleutelvereiste voor effectieve governance. In het kader van deze bachelorproef wordt transparantie niet opgevat als publieke openbaarheid van gegevens, maar als interne, governancegerichte transparantie binnen VLAIO. Dit betekent dat data engineers, data analisten, auditors en rapportgebruikers inzicht hebben in de oorsprong, transformaties en afhankelijkheden van data assets. Deze vorm van transparantie ondersteunt de interne controleerbaarheid en verantwoordingsplicht, blijvend binnen de voorwaarden van GDPR, zoals dataminimalisatie\footnote{Persoonsgegevens mogen enkel verzameld en verwerkt worden indien deze noodzakelijk zijn voor het specifieke doel, waarvoor ze worden gebruikt.} en gepaste toegangscontrole. Daarnaast zorgt deze herleidbaarheid ook voor governance doelen zoals verantwoordingsplicht (accountability), kwaliteitsgarantie, informatiebeveiliging en impactanalyse \autocite{verma2024tracing}.\\

Naast deze governance doelen, speelt lineage ook een belangrijke rol als ondersteuning binnen regelgevings- en privacycontexten zoals GDPR (General Data Protection Regulation) of andere compliancevereisten. \textcite{malviya2025dbt} tonen aan dat een transparante data lineage controleerbaarheid (auditability) sterk verhoogt. Zo worden DSAR-verzoeken\footnote{Verzoeken van betrokkenen in het kader van GDPR, bijvoorbeeld voor inzage, correctie of verwijdering van persoonsgegevens.} 64\% sneller afgehandeld en audit-evidence\footnote{Bewijsmateriaal dat aantoont dat processen, controles en gegevensverwerkingen voldoen aan geldende regelgeving en interne governance-eisen} 80\% efficiënter samengesteld. Hierdoor kan een organisatie beter demonstreren dat ze regelgeving als GDPR naleven. Bovendien kan de transparantie ervoor zorgen dat een organisatie een betere relatie kan opbouwen met zijn stakeholders en met meer vertrouwen belangrijke beslissingen nemen \autocite{verma2024tracing}. Deze voordelen zijn zeker interessant, maar niet van toepassing in dit onderzoek.\\

In deze bachelorproef wordt dieper ingegaan op impactanalyse. Meer specifiek wordt onderzocht hoe data lineage als governance-tool kan bijdragen aan de transparantie, herleidbaarheid en betrouwbaarheid van impactanalyses.\\

\section{Data lineage en impactanalyse}

Change Impact Analysis ofwel impactanalyse wordt door \textcite{kretsou2021change} omschreven als het proces waarbij wordt nagegaan welke gevolgen een wijziging kan hebben voor andere onderdelen van een systeem. Binnen een data omgeving wil dit concreet zeggen dat wijzigingen binnen datasets, datamodellen of transformatieprocessen worden geëvalueerd op de downstream-effecten die ze hebben. Binnen de context van databronnen wordt er een onderscheid gemaakt tussen upstream en downstream bronnen. Zo verwijst upstream naar databronnen of processen die data leveren aan een bepaald punt in de dataketen en downstream verwijst naar systemen of processen die de data ontvangen of gebruiken. Zonder inzicht in de onderlinge afhankelijkheden tussen databronnen, transformaties en rapportages wordt het echter moeilijk om te bepalen welke onderdelen van een organisatie beïnvloed worden, wanneer een bepaalde wijziging doorgevoerd wordt.\\

Met lineage kan je duidelijk zien welke upstream-bronnen bijdragen aan specifieke datasets en welke downstream rapporten, dashboards of andere processen van deze data afhankelijk zijn. Dit kan waardevol zijn bij het uitvoeren van betrouwbare impactanalyses. Data lineage maakt het mogelijk om veranderingen gericht te analyseren en evalueren voordat ze in de praktijk worden doorgevoerd. Zonder dit inzicht wordt het achterhalen van de impact die wijzigingen met zich mee brengen vaak een manueel en tijdrovend proces, aangezien je elke afhankelijkheid per dataset of zelfs per veld moet controleren \autocite{malviya2025dbt}. Wanneer er binnen een organisatie een voorstel wordt gedaan voor een wijziging, als een systeemupgrade of een data transformatie, kan je dus aan de hand van data lineage zien welke downstream processen of data gebruikers er geraakt zullen worden. Dit zorgt vervolgens voor een betere risicobeoordeling en mitigatieplanning, het voorkomen van negatieve gevolgen van een specifiek risico of proces \autocite{verma2024tracing}.\\

Daarnaast ondersteunen data lineage en provenance de reproduceerbaarheid van resultaten doordat zij traceerbaarheid mogelijk maken. Het documenteren van lineage en verwerkingsstappen maakt het haalbaar om inputs en outputs van processen in een gedistribueerde dataketen terug te vinden \autocite{wittnerlightweight}. Dat verhoogt de betrouwbaarheid van de analyses en de mogelijkheid om eerder gemaakte keuzes en resultaten te reconstrueren \autocite{verma2024tracing}.\\

Het is belangrijk te benadrukken dat de inhoud en context van impactanalyses organisatieafhankelijk zijn. Bij VLAIO is een impactanalyse niet enkel een technische oefening waarbij afhankelijkheden tussen datasets worden opgespoord, maar omvat ook het toelichten van de impact van wijzigingen aan verschillende betrokkenen, het documenteren van beslissingen en het ondersteunen van de verdere uitvoering in rapportering en datapijplijnen. Hoewel impactanalyse in de literatuur vaak wordt benaderd als een technisch proces waarbij de gevolgen van wijzigingen in datasets, transformaties of rapportagesystemen systematisch in kaart worden gebracht, wordt het concept in deze bachelorproef ruimer benaderd binnen een data governance-context. Impactanalyse wordt daarbij niet alleen opgevat als een formele activiteit voor data engineers of ontwikkelaars, maar ook als een proces waarbij andere gebruikers van data, zoals analisten en businessgebruikers, nood hebben aan inzicht in de herkomst, betekenis en betrouwbaarheid van gegevens. Wanneer zij geconfronteerd worden met gewijzigde rapporten, onverwachte waarden of moeilijk verklaarbare verschillen in output, stellen ook zij in essentie vragen naar de impact van onderliggende wijzigingen. In dat opzicht verschuift de focus van een technische afhankelijkheidsdetectie naar vragen rond vertrouwen, betekenis en herkomst van data. Data lineage kan hierin een belangrijke rol spelen, doordat het niet alleen formele impactanalyses ondersteunt, maar ook bredere transparantie en herleidbaarheid mogelijk maakt voor een ruimer gebruikerspubliek. De waargenomen verbeteringen zijn daarbij niet uitsluitend technisch, maar hebben ook te maken met het beter vastleggen en toegankelijk maken van kennis die eerder vooral impliciet of verspreid aanwezig was. Deze bredere interpretatie sluit aan bij het governance-perspectief van deze bachelorproef, waarin de meerwaarde van lineage zowel technisch als vanuit begrijpelijkheid en bruikbaarheid voor stakeholders wordt benaderd.

\section{Meetbaarheid van data governance en procesverbetering}

In de literatuur worden de strategische en operationele voordelen van data lineage binnen data governance en impactanalyse duidelijk naar voren gebracht \autocite{bernardo2024data,verma2024tracing}. Het aantonen van die voordelen in de praktijk blijkt echter minder vanzelfsprekend. Begrippen als transparantie, verantwoordingsplicht en vertrouwen worden frequent genoemd, maar worden zelden vertaald naar hoe gebruikers deze voordelen in hun dagelijkse werk ervaren, bijvoorbeeld in termen van begrijpbaarheid, bruikbaarheid of kwaliteit van rapporten.\\

Wanneer de focus ligt op impactanalyse in een concrete organisatiecontext, kan de bijdrage van lineage daarom ook op een praktijkgerichte manier geëvalueerd worden. Naast mogelijke observaties van het analyseproces kan daarbij vooral gekeken worden naar hoe gebruikers de beschikbaarheid en bruikbaarheid van lineage-informatie ervaren: begrijpen zij beter waar data vandaan komt, hebben zij meer vertrouwen in de rapportering en kunnen zij afhankelijkheden duidelijker reconstrueren en communiceren?\\

Een dergelijke evaluatie is dus zeker relevant omdat de meerwaarde van lineage niet uitsluitend in technische verbeteringen ligt. Verbeteringen hebben ook te maken met hoe kennis binnen de organisatie gedeeld en gebruikt wordt. Zo wordt informatie die vroeger vooral impliciet of verspreid aanwezig was, beter beschreven en vindbaar. De evaluatie van een lineage-oplossing moet daarom niet alleen letten op de technische zichtbaarheid van afhankelijkheden, maar ook op begrijpelijkheid, bruikbaarheid en de ervaren ondersteuning bij impactanalyses..\\

\section{Metadata management tools en data catalogs}

In essentie steunt data lineage op metadata. Metadata is beschrijvende informatie over data zoals onder andere welke data assets er bestaan, semantiek van data, waar deze assets zich bevinden, door welke systemen en transformaties ze bewegen en of deze confidentieel zijn \autocite{eichler2021enterprise}. Metadata management is een verzameling van alle activiteiten om data assets via hun metadata te beheren. Zonder metadata weet een organisatie niet welke data ze bezitten, wat die representeert en of ze vertrouwelijk is, waardoor doelstellingen als compliance en het volledig kunnen benutten van datawaarde in bedwang komen \autocite{eichler2021enterprise}. Volgens \textcite{eichler2021enterprise} is datatransparantie het kunnen vinden, begrijpen en raadplegen van data, wat door metadata management kan worden gewaarborgd.\\

Een data catalog is de technologische tool voor metadata management. Het werkt als een centraal data inventaris met bijhorende documentatie over de geregistreerde databronnen, business, technische en operationele metadata \autocite{eichler2021enterprise}. Business metadata omvat bijvoorbeeld een definitie of inhoudelijke beschrijving van "customer purchase", technische metadata kan bijvoorbeeld datatype informatie bevatten en operationele metadata beschrijft bijvoorbeeld de gebruiksgeschiedenis van data. Door al deze metadata te centraliseren en beschikbaar te maken in één interface, ondersteunen data catalogs de begrijpbaarheid en vindbaarheid van data over meerdere systemen heen, zowel voor IT als Business.\\

De meeste moderne metadata platformen gaan echter verder dan enkel data in\-ven\-tarisatie. Lineage functionaliteit in tools steunt typisch op technische metadata uit systemen zoals bijvoorbeeld schema's, ETL-workflows en logbestanden om datastromen te reconstrueren en te visualiseren, waardoor upstream- en downstreamafhankelijkheden expliciet in kaart worden gebracht en impactanalyses bij wijzigingen rechtstreeks worden ondersteund \autocite{verma2024tracing}. In complexe en gedistribueerde data architecturen zoals datalakes en ERP structuren is een organisatiebrede aanpak zeker interessant, aangezien het voor een individuele medewerker onhaalbaar wordt om een volledig overzicht van alle data en systemen bij te houden \autocite{eichler2021enterprise}. Dergelijke platformen zijn vooral relevant wanneer afhankelijkheden verspreid zitten over verschillende databronnen, transformatielagen en rapporteringstools, waardoor manuele reconstructie van de volledige dataflow moeilijk is en zeer veel tijd inneemt.\\

Binnen dit oplossingsdomein bestaan er zowel commerciële als open source platformen die data cataloging, governance en lineage functionaliteit combineren \autocite{eichler2021enterprise}. In deze bachelorproef dient OpenMetadata als representatief voorbeeld van een open source metadata en governanceplatform met geïntegreerde lineage-functionaliteit.\\

\section{Positionering van OpenMetadata}

OpenMetadata wordt in deze bachelorproef gepositioneerd als een open source metadataplatform dat functies biedt voor zowel data cataloging, data governance en data lineage. Het platform omschrijft zichzelf als een geïntegreerd metadata platform ofwel een plek waar alle metadata van een organisatie overzichtelijk samenkomt. Een centrale plek dus voor data discovery, observability en governance, ondersteund door een metadata repository en lineage \autocite{OpenMetadataDocs2024}. Binnen VLAIO werd OpenMetadata gekozen als metadata-managementtool na een evaluatie naast 17 alternatieve tools, waaronder bijvoorbeeld Microsoft Purview. De betrokkenen gaven aan dat Purview steeds minder bruikbaar werd voor VLAIO. Het product is verschoven van data governance naar security en DLP. De nadruk ligt nu vooral op beveiligingsrapporten, terwijl VLAIO juist behoefte heeft aan data lineage, duidelijke definities, ownership en ondersteuning bij impactanalyses. OpenMetadata werd dus beoordeeld als beter afgestemd op de governance-doelstellingen van VLAIO. Daarnaast is het ook open-source, waardoor er geen sprake is van vendor lock-in.\\

OpenMetadata past binnen het oplossingsdomein van data catalogs en metadata management omdat de catalogfunctie metadata rond data assets centraal beschikbaar maakt voor vindbaarheid en begrip \autocite{eichler2021enterprise}. De lineage functionaliteit maakt bovendien upstream en downstream afhankelijkheden zichtbaar, wat rechtstreeks kan bijdragen aan transparantere en beter onderbouwde impactanalyses bij wijzigingen \autocite{verma2024tracing,OpenMetadataDocs2024}. In deze bachelorproef wordt binnen de context van VLAIO nagegaan in welke mate OpenMetadata kan bijdragen aan transparantere en beter herleidbare impactanalyses.\\


\section{Onderzoeksleemte}

Hoewel literatuur duidelijk aantoont dat data lineage bijdraagt aan transparantie en impactanalyse binnen data governance, blijft het in veel gevallen onduidelijk hoe deze meerwaarde zich manifesteert in een concrete organisatiecontext \autocite{verma2024tracing,bernardo2024data}. Bestaande studies bespreken lineage vaak op conceptueel of technisch niveau, maar besteden minder aandacht aan hoe gebruikers de beschikbaarheid van lineage-informatie nu ervaren. Omdat organisatiebreed metadatamanagement en de keuze voor de juiste tools in de praktijk vaak ingewikkeld blijken, is er een sterke behoefte aan praktijkgerichte evaluaties \autocite{eichler2021enterprise}. In deze bachelorproef wordt daarom binnen de context van VLAIO onderzocht hoe OpenMetadata de transparantie en herleidbaarheid van impactanalyses en rapporten kan ondersteunen, met bijzondere aandacht voor begrijpelijkheid, bruikbaarheid en ervaren meerwaarde voor betrokken stakeholders.\\